%% ============================================================================
%% Fichier  : annexes/Ann_D_lois.tex
%% Titre    : Textes juridiques et normatifs
%% Version  : 1.0 -- novembre 2026
%% ============================================================================
%% Alignée sur les annexes des autres cours de l'auteur (ACSI, INF3226).
%% ============================================================================

\chapter{Textes Juridiques et Normatifs}
\label{ann:lois}

Cette annexe rassemble les textes cités dans le cours. Seuls les extraits lus
dans leur version officielle sont reproduits entre guillemets~; les autres sont
résumés, avec leur référence.

\section*{Loi n° 2010/012 relative à la cybersécurité et à la cybercriminalité}

\begin{boxjuris}
\textbf{Texte officiel}~\cite{cameroun2010a}

\medskip
\textbf{Article~25, alinéas~1 et~2}~: \textit{« (1) Les opérateurs de réseaux et les fournisseurs de services de communications électroniques ont l'obligation de conserver les données de connexion et de trafic pendant une période de dix (10) ans. (2) Les opérateurs de réseaux et les fournisseurs de services de communications électroniques installent des mécanismes de surveillance de trafic des données de leurs réseaux. Ces données peuvent être accessibles lors des investigations judiciaires. »}

\medskip
\textbf{Article~52, alinéa~3}~: \textit{« Les Officiers de Police Judiciaire
et les agents habilités de l'Agence peuvent, lors des investigations,
accéder aux moyens de transport, à tout local à usage professionnel, à
l'exclusion des domiciles privés, en vue de rechercher, de constater les
infractions, de demander la communication de tous les documents
professionnels et en prendre copie, recueillir, sur convocation ou sur
place, les renseignements et justifications. »}

\medskip
\textbf{Article~53, alinéa~1}~: \textit{« Les perquisitions en matière de
cybercriminalité sont susceptibles de porter sur des données qui peuvent
être des supports physiques ou des copies réalisées en présence des
personnes qui assistent à la perquisition. »}

\medskip
\textbf{Article~54}~: \textit{« Les perquisitions et les saisies sont effectuées conformément aux dispositions du Code de Procédure Pénale en tenant compte du dépérissement des preuves. »}

\medskip
\textbf{Article~55}~: \textit{« (1) Lorsqu'il apparaît que les données saisies ou obtenues au cours de l'enquête ou de l'instruction ont fait l'objet d'opérations de transformation empêchant d'accéder en clair ou sont de nature à compromettre les informations qu'elles contiennent, le Procureur de la République, le Juge d'Instruction ou la juridiction de jugement peuvent réquisitionner toute personne physique ou morale qualifiée, en vue d'effectuer les opérations techniques permettant d'obtenir la version en clair desdites données. (2) Lorsqu'un moyen de cryptographie a été utilisé, les autorités judiciaires peuvent exiger la convention secrète de déchiffrement du cryptogramme. »}

\medskip
\textbf{Article~65, alinéa~2}~: \textit{« Est puni des peines prévues à
l'alinéa 1 ci-dessus, tout accès non autorisé, à l'ensemble ou à une
partie d'un réseau de communications électroniques ou d'un système
d'information ou d'un équipement terminal. »}

\medskip
\textbf{Article~69}~: \textit{« Est puni d'un emprisonnement de cinq (05) à
dix (10) ans et d'une amende de 10.000.000 (dix millions) à 100.000.000
(cent millions) F CFA ou de l'une de ces deux peines seulement, celui qui
accède sans droit, et en violation des mesures de sécurité, à l'ensemble ou
à une partie d'un réseau de communications électroniques, d'un système
d'information ou d'un équipement terminal, afin d'obtenir des informations
ou des données, en relation avec un système d'information connecté à un
autre système d'information. »}

\medskip
\textbf{Article~70}~: \textit{« Est puni d'une amende de 1.000.000 (un
million) à 5.000.000 (cinq millions) F CFA, celui qui provoque par
saturation, l'attaque d'une ressource de réseau de communications
électroniques ou d'un système d'information dans le but de l'effondrer en
empêchant la réalisation des services attendus. »}

\medskip
\textbf{Article~86, alinéa~1}~: \textit{« Est puni des peines prévues
l'article 71 ci-dessus, celui qui importe, détient, offre, cède, vend ou
met à disposition, sous quelle que forme que ce soit, un programme
informatique, un mot de passe, un code d'accès ou toutes données
informatiques similaires conçus et ou spécialement adaptés, pour permettre
d'accéder, à tout ou partie d'un réseau de communications électroniques ou
d'un système d'information. »}
\end{boxjuris}

\section*{Autres textes}

\begin{tableaucours}{L{4.4cm}L{4.6cm}L{3.2cm}}{%
    \tableauHeader{Texte} &
    \tableauHeader{Objet} &
    \tableauHeader{Vérification}
}
Loi n°~2010/013~\cite{cameroun2010b} & Communications électroniques & Présentée~; articles à vérifier \\
Loi n°~2024/017~\cite{cameroun2024donnees} & Protection des données personnelles & Présentée~; texte à collationner au Journal officiel \\
Convention de Budapest~\cite{budapest2001} & Cybercriminalité, entraide, conservation rapide & Texte du Conseil de l'Europe \\
Convention de Malabo~\cite{malabo2014} & Cybersécurité et données personnelles (UA) & Texte de l'Union africaine \\
RGPD~\cite{rgpd2016} & Données personnelles (UE) & Journal officiel de l'UE \\
CLOUD Act~\cite{cloudact2018} & Accès aux données détenues par des fournisseurs américains & Texte fédéral des États-Unis \\
\end{tableaucours}

\section*{Normes de la preuve numérique}

\begin{tableaucours}{L{3.6cm}L{8.8cm}}{%
    \tableauHeader{Norme} &
    \tableauHeader{Objet}
}
ISO/IEC~27037~\cite{iso27037} & Identification, collecte, acquisition et préservation \\
ISO/IEC~27041~\cite{iso27041_2015} & Assurer l'aptitude des méthodes d'investigation \\
ISO/IEC~27042~\cite{iso27042_2015} & Analyse et interprétation de la preuve numérique \\
ISO/IEC~27043~\cite{iso27043_2015} & Principes et processus d'investigation \\
NIST SP~800-86~\cite{nist80086} & Intégrer la forensique à la réponse aux incidents \\
RFC~3227~\cite{rfc3227} & Collecte et archivage des preuves, ordre de volatilité \\
Guide de l'ACPO~\cite{acpo2012} & Quatre principes de la preuve numérique \\
\end{tableaucours}
